\heading{实验物理中的统计方法-23春-期末}
\noteinfo{内容精修版；仅保留正文内容，适配已有 preamble 与题目环境。}
\begin{question}[误差传播]
两个物理量 $X$ 和 $Y$ 的标准不确定度分别为 $\sigma_X,\sigma_Y$。求 $X-Y$ 的标准不确定度。分别给出 $X,Y$ 独立和不独立时的公式。
\begin{solution}
若 $X,Y$ 独立，则
\[
\V(X-Y)=\V(X)+\V(Y),
\]
所以
\ansmath{\sigma_{X-Y}=\sqrt{\sigma_X^2+\sigma_Y^2}}.
若 $X,Y$ 不独立，则要保留协方差项：
\[
\V(X-Y)=\V(X)+\V(Y)-2\cov(X,Y),
\]
即
\ansmath{\sigma_{X-Y}=\sqrt{\sigma_X^2+\sigma_Y^2-2\cov(X,Y)}}.
若记相关系数为 $\rho$，则也可写为
\[
\sigma_{X-Y}=\sqrt{\sigma_X^2+\sigma_Y^2-2\rho\sigma_X\sigma_Y}.
\]
\end{solution}
\end{question}

\begin{question}[卡方分布的渐近正态]
设 $X\sim\chi^2(2n)$。当 $n$ 很大时，
\[
\frac{X-2n}{2\sqrt n}
\]
近似服从什么分布？
\begin{solution}
若 $X\sim\chi^2(k)$，则
\[
\E[X]=k,\qquad \V(X)=2k.
\]
这里 $k=2n$，故
\[
\E[X]=2n,\qquad \V(X)=4n.
\]
卡方分布是独立标准正态平方和。当自由度趋于无穷时，经标准化后满足渐近正态性，因此
\ansmath{\frac{X-2n}{2\sqrt n}\xrightarrow[]{d} \N(0,1)}.
\end{solution}
\end{question}

\begin{question}[正态总体的常用抽样分布]
设 $X_1,\ldots,X_n$ 是来自正态总体 $\N(\mu,\sigma^2)$ 的简单随机样本，定义
\[
\bar X=\frac1n\sum_{i=1}^nX_i,\qquad
S^2=\frac1{n-1}\sum_{i=1}^n(X_i-\bar X)^2.
\]
写出 $\bar X$、$(n-1)S^2/\sigma^2$ 以及 $\sqrt n(\bar X-\mu)/S$ 的分布。
\begin{solution}
正态总体样本均值仍服从正态分布：
\[
\bar X\sim \N\left(\mu,\frac{\sigma^2}{n}\right).
\]
样本方差满足
\[
\frac{(n-1)S^2}{\sigma^2}\sim\chi^2(n-1).
\]
并且在正态总体下，$\bar X$ 与 $S^2$ 独立，所以
\[
\frac{\sqrt n(\bar X-\mu)}{S}
=\frac{(\bar X-\mu)/(\sigma/\sqrt n)}{\sqrt{[(n-1)S^2/\sigma^2]/(n-1)}}
\sim t(n-1).
\]
\ansmath{\bar X\sim \N\left(\mu,\frac{\sigma^2}{n}\right),\quad
\frac{(n-1)S^2}{\sigma^2}\sim\chi^2(n-1),\quad
\frac{\sqrt n(\bar X-\mu)}{S}\sim t(n-1)}
\end{solution}
\end{question}

\begin{question}[矩估计与极大似然估计]
设随机变量 $X$ 的概率密度为
\[
f(x;\theta)=
\begin{cases}
(\theta+1)x^\theta,&0<x<1,\\
0,&\text{其他},
\end{cases}
\qquad \theta>-1.
\]
现有独立样本 $X_1,\ldots,X_n$，求 $\theta$ 的矩估计量和极大似然估计量。
\begin{solution}
先求一阶矩：
\[
\E[X]=\int_0^1 x(\theta+1)x^\theta\,\dd x
=(\theta+1)\int_0^1x^{\theta+1}\,\dd x
=\frac{\theta+1}{\theta+2}.
\]
令 $\bar X=(\theta+1)/(\theta+2)$，解得
\[
\bar X(\theta+2)=\theta+1,
\qquad
\theta(\bar X-1)=1-2\bar X,
\]
所以矩估计量为
\ansmath{\hat\theta_{\rm MOM}=\frac{2\bar X-1}{1-\bar X}}.

似然函数为
\[
L(\theta)=\prod_{i=1}^n(\theta+1)X_i^\theta
=(\theta+1)^n\left(\prod_{i=1}^nX_i\right)^\theta,
\qquad \theta>-1.
\]
对数似然为
\[
\ell(\theta)=n\ln(\theta+1)+\theta\sum_{i=1}^n\ln X_i.
\]
求导：
\[
\ell'(\theta)=\frac{n}{\theta+1}+\sum_{i=1}^n\ln X_i,
\qquad
\ell''(\theta)=-\frac{n}{(\theta+1)^2}<0.
\]
令 $\ell'(\theta)=0$ 得
\ansmath{\hat\theta_{\rm MLE}=-\frac{n}{\sum_{i=1}^n\ln X_i}-1}.
由于 $0<X_i<1$ 几乎必然成立，故 $\sum_i\ln X_i<0$，该估计量满足 $\hat\theta_{\rm MLE}>-1$。
\end{solution}
\end{question}

\begin{question}[假设检验错误]
在假设检验中，若原假设 $H_0$ 为真却被拒绝，这类错误叫做\underline{\hspace{2.5em}}；若 $H_0$ 为假却没有被拒绝，这类错误叫做\underline{\hspace{2.5em}}。怎样才能同时减少这两类错误的概率？
\begin{solution}
$H_0$ 为真却被拒绝称为第一类错误，也叫弃真错误；$H_0$ 为假却没有被拒绝称为第二类错误，也叫存伪错误或取伪错误。
\anstext{第一类错误；第二类错误}

需要注意，“接受 $H_0$”的表述不如“没有拒绝 $H_0$”严谨，因为检验通常只能说明证据不足以拒绝 $H_0$。

在样本量和检验形式固定时，降低显著性水平 $\alpha$ 往往会增大第二类错误概率 $\beta$，二者通常不能任意同时降低。若要同时降低两类错误概率，常用办法是增加样本量、降低测量不确定度，或针对明确备择假设采用功效更高的检验方法。
\end{solution}
\end{question}

\begin{question}[泊松均值的有效估计]
设一次计数实验的观测值为 $N=n$，并假定 $N\sim\Pois(\lambda)$。求 $\lambda$ 的极大似然估计量，并证明它是无偏且有效的估计量。
\begin{solution}
似然函数为
\[
L(\lambda)=P(N=n\mid \lambda)=\ee^{-\lambda}\frac{\lambda^n}{n!},
\qquad \lambda>0.
\]
对数似然为
\[
\ell(\lambda)=-\lambda+n\ln\lambda-\ln n!.
\]
求导得
\[
\ell'(\lambda)=-1+\frac{n}{\lambda},
\]
令其为零，得到
\[
\hat\lambda_{\rm MLE}=n.
\]
作为随机变量，估计量就是 $\hat\lambda=N$。

由于泊松分布满足
\[
\E[N]=\lambda,
\qquad
\V(N)=\lambda,
\]
所以 $\hat\lambda=N$ 是无偏估计量，且方差为 $\lambda$。

单个观测的 Fisher 信息为
\[
I(\lambda)=\E\left[\left(\frac{\partial}{\partial\lambda}\ln P(N\mid\lambda)\right)^2\right]
=\E\left[\left(-1+\frac{N}{\lambda}\right)^2\right]
=\frac{\V(N)}{\lambda^2}=\frac1\lambda.
\]
Cram\'er--Rao 下界为
\[
\V(\hat\lambda)\ge \frac1{I(\lambda)}=\lambda.
\]
而 $\V(N)=\lambda$ 正好达到下界，因此
\anstext{$\hat\lambda=N$ 是有效估计量。}
\end{solution}
\end{question}

\begin{question}[单样本 $t$ 检验]
从某班抽出 $25$ 名学生，数学成绩样本均值为 $84$ 分，样本标准差为 $8$ 分。假定该班成绩可看作来自正态总体，且总体方差未知。问在 $90\%$ 置信水平下，是否可以认为该班数学成绩总体均值为年段均值 $87$ 分？已知 $t_{0.95}(24)=1.711$。
\begin{solution}
这是总体方差未知的单样本双侧 $t$ 检验。设
\[
H_0:\mu=87,
\qquad
H_1:\mu\ne87.
\]
检验统计量为
\[
T=\frac{\bar X-\mu_0}{S/\sqrt n}
=\frac{84-87}{8/\sqrt{25}}
=\frac{-3}{1.6}=-1.875.
\]
$90\%$ 置信水平对应双侧显著性水平 $\alpha=0.10$，临界值为
\[
t_{1-\alpha/2}(24)=t_{0.95}(24)=1.711.
\]
因为
\[
|T|=1.875>1.711,
\]
故在 $\alpha=0.10$ 水平下拒绝 $H_0$。因此不能认为该班总体均值等于 $87$ 分。

等价地，$90\%$ 置信区间为
\[
\bar X\pm t_{0.95}(24)\frac{S}{\sqrt n}
=84\pm1.711\times1.6
=[81.26,86.74],
\]
其中不包含 $87$。所以结论一致。
\anstext{拒绝 $\mu=87$；在 $90\%$ 置信水平下不能认为该班均值为 $87$ 分。}
\end{solution}
\end{question}

\begin{question}[均匀分布端点估计]
设 $X_1,\ldots,X_n$ 是来自均匀分布 $U(0,\theta)$ 的独立样本，其中 $\theta>0$。记
\[
M=X_{(n)}=\max\{X_1,\ldots,X_n\}.
\]
\begin{enumerate}
\item 求 $\theta$ 的极大似然估计 $\hat\theta$，并求 $Y=\hat\theta/\theta$ 的概率密度函数。
\item 将 $\hat\theta$ 乘以常数 $c$ 构造无偏估计量，求 $c$。按均方误差比较，修正后的估计量和原极大似然估计量哪个更好？
\item 利用 $Y$ 给出 $\theta$ 的置信水平为 $1-\alpha$ 的双侧等尾置信区间。若 $n=10$ 且 $M=9.8$，给出 $95\%$ 置信区间。
\end{enumerate}
\begin{solution}
\begin{enumerate}
\item 对给定样本，似然函数为
\[
L(\theta)=\theta^{-n}\mathbf 1_{\{\theta\ge M\}}.
\]
在 $\theta\ge M$ 上，$\theta^{-n}$ 随 $\theta$ 单调递减，因此极大似然估计为
\ansmath{\hat\theta=M}.

先求 $M$ 的分布：当 $0<m<\theta$ 时，
\[
P(M\le m)=P(X_1\le m,\ldots,X_n\le m)=\left(\frac{m}{\theta}\right)^n.
\]
令 $Y=M/\theta$，则 $0<Y<1$，且
\[
P(Y\le y)=P(M\le \theta y)=y^n,
\qquad 0<y<1.
\]
所以
\ansmath{g_Y(y)=ny^{n-1},\qquad 0<y<1}.

\item 由 $Y=M/\theta$ 的密度可得
\[
\E[Y]=\int_0^1y\,ny^{n-1}\,\dd y=\frac{n}{n+1}.
\]
因此
\[
\E[M]=\frac{n}{n+1}\theta.
\]
要使 $cM$ 无偏，需要
\[
c\frac{n}{n+1}\theta=\theta,
\]
故
\ansmath{c=\frac{n+1}{n},\qquad \tilde\theta=\frac{n+1}{n}M}.

进一步，
\[
\E[Y^2]=\int_0^1 y^2ny^{n-1}\,\dd y=\frac{n}{n+2},
\]
所以
\[
\V(M)=\theta^2\left(\frac{n}{n+2}-\frac{n^2}{(n+1)^2}\right)
=\frac{n\theta^2}{(n+1)^2(n+2)}.
\]
无偏修正估计量的方差为
\[
\V(\tilde\theta)=\left(\frac{n+1}{n}\right)^2\V(M)
=\frac{\theta^2}{n(n+2)}.
\]
原 MLE 的偏差为
\[
\operatorname{Bias}(M)=\E[M]-\theta=-\frac{\theta}{n+1},
\]
故
\[
\operatorname{MSE}(M)=\V(M)+\operatorname{Bias}(M)^2
=\frac{2\theta^2}{(n+1)(n+2)}.
\]
比较
\[
\frac{\theta^2}{n(n+2)}<\frac{2\theta^2}{(n+1)(n+2)}
\quad\Longleftrightarrow\quad n>1.
\]
因此当 $n>1$ 时，按均方误差比较，无偏修正估计量 $\tilde\theta$ 更好；当 $n=1$ 时二者均方误差相同。

\item 因 $Y=M/\theta$ 的分布与 $\theta$ 无关，取
\[
P\left(a\le \frac{M}{\theta}\le b\right)=1-\alpha.
\]
等尾选择满足
\[
P(Y<a)=\frac{\alpha}{2},\qquad P(Y>b)=\frac{\alpha}{2}.
\]
由于 $F_Y(y)=y^n$，
\[
a=\left(\frac{\alpha}{2}\right)^{1/n},
\qquad
b=\left(1-\frac{\alpha}{2}\right)^{1/n}.
\]
由
\[
a\le \frac{M}{\theta}\le b
\]
反解得
\ansmath{\theta\in\left[\frac{M}{b},\frac{M}{a}\right]
=\left[\frac{M}{(1-\alpha/2)^{1/n}},\frac{M}{(\alpha/2)^{1/n}}\right]}.

当 $n=10$、$M=9.8$、$\alpha=0.05$ 时，
\[
a=0.025^{1/10}\approx0.6915,
\qquad
b=0.975^{1/10}\approx0.9975.
\]
所以
\[
\theta\in\left[\frac{9.8}{0.9975},\frac{9.8}{0.6915}\right]
\approx[9.825,14.172].
\]
\anstext{$95\%$ 等尾置信区间约为 $[9.83,14.17]$。}
\end{enumerate}
\end{solution}
\end{question}

\begin{question}[过原点的加权最小二乘]
小车做匀速直线运动。已知小车在 $t=0$ 时通过 $d=0$，但这一点不作为实验数据。现测得 6 组 $(d,t)$ 数据如下，距离 $d$ 视为准确，时间测量的标准不确定度为 $\sigma_t=0.1\,\mathrm{s}$：
\[
\begin{array}{c|cccccc}
 d\ (\mathrm m) &0.50&1.00&1.50&2.00&2.50&3.00\\ \hline
 t\ (\mathrm s) &1.02&2.07&2.94&4.11&4.92&6.06
\end{array}
\]
用过原点的最小二乘法估计小车速度 $v$，并给出 $\chi^2_{\min}$。同时写出等方差情形的一般公式。
\begin{solution}
因为主要不确定度在 $t$ 上，应将模型写为
\[
t_i=kd_i+\epsilon_i,
\qquad k=\frac1v,
\qquad \epsilon_i\sim\N(0,\sigma_t^2).
\]
由于原点不是实验数据，但运动模型经过原点，所以不拟合截距。于是
\[
\chi^2(k)=\sum_{i=1}^6\frac{(t_i-kd_i)^2}{\sigma_t^2}.
\]
令导数为零：
\[
\frac{\dd}{\dd k}\sum_{i=1}^6(t_i-kd_i)^2
=-2\sum_{i=1}^6d_i(t_i-kd_i)=0,
\]
得
\ansmath{\hat k=\frac{\sum_i d_it_i}{\sum_i d_i^2},\qquad
\hat v=\frac1{\hat k}=\frac{\sum_i d_i^2}{\sum_i d_it_i}}.
对本题数据，
\[
\sum_i d_i^2=22.75,
\qquad
\sum_i d_it_i=45.69,
\]
所以
\[
\hat k=\frac{45.69}{22.75}=2.00835\ \mathrm{s/m},
\qquad
\hat v=0.49792\ \mathrm{m/s}.
\]
代回可得
\[
\chi^2_{\min}
=\frac1{\sigma_t^2}\left(\sum_i t_i^2-\frac{(\sum_i d_it_i)^2}{\sum_i d_i^2}\right).
\]
本题中
\[
\sum_i t_i^2=91.791,
\]
因此
\[
\chi^2_{\min}
=\frac1{0.1^2}\left(91.791-\frac{45.69^2}{22.75}\right)
\approx2.94.
\]
\anstext{$\hat v\approx0.498\ \mathrm{m/s}$，$\chi^2_{\min}\approx2.94$。}
若需要速度的不确定度，在 $\sigma_t$ 已知时有
\[
\sigma_{\hat k}=\frac{\sigma_t}{\sqrt{\sum_i d_i^2}},
\qquad
\sigma_{\hat v}\approx\frac{\sigma_{\hat k}}{\hat k^2},
\]
本题约为 $\sigma_{\hat v}=0.0052\,\mathrm{m/s}$。
\end{solution}
\end{question}

\begin{question}[泊松计数与寿命置信区间]
某探测器寻找质子衰变。事例数可认为服从泊松分布，不考虑本底。设探测器中可衰变质子数为
\[
N_p=3.35\times10^{29},
\]
观测时间为 $T=1$ 年，实际观察到 $n=7$ 个候选事例。给出质子寿命 $\tau$ 的双侧等尾 $90\%$ 置信区间。已知泊松均值 $\mu$ 的等尾置信区间可写为
\[
\mu\in\left[\frac12\chi^2_{2n,\alpha/2},\frac12\chi^2_{2n+2,1-\alpha/2}\right].
\]
\begin{solution}
若质子寿命为 $\tau$，则期望衰变事例数为
\[
\mu=\frac{N_pT}{\tau}.
\]
本题 $n=7$，置信水平 $90\%$ 对应 $\alpha=0.10$。泊松均值的等尾区间为
\[
\mu\in\left[\frac12\chi^2_{14,0.05},\frac12\chi^2_{16,0.95}\right].
\]
使用分位数
\[
\chi^2_{14,0.05}=6.571,
\qquad
\chi^2_{16,0.95}=26.296,
\]
得到
\[
\mu\in[3.2855,13.148].
\]
由于
\[
\tau=\frac{N_pT}{\mu},
\]
且 $\tau$ 与 $\mu$ 单调反向，对应的寿命区间为
\[
\tau\in\left[\frac{N_pT}{13.148},\frac{N_pT}{3.2855}\right].
\]
代入 $N_p=3.35\times10^{29}$、$T=1$ 年，得
\[
\tau\in[2.55\times10^{28},1.02\times10^{29}]\ \text{年}.
\]
\anstext{$90\%$ 等尾置信区间约为 $[2.55\times10^{28},1.02\times10^{29}]$ 年。}
\end{solution}
\end{question}

\begin{question}[相关系统误差下的最佳平均]
同一个物理量 $x$ 用同一台仪器或同一外部输入测量两次，得到
\[
x_1\pm\sigma_{s1}\pm\sigma_{c1},
\qquad
x_2\pm\sigma_{s2}\pm\sigma_{c2}.
\]
其中第一个不确定度为统计不确定度，第二个为由共同来源引入的系统不确定度。设两次统计误差相互独立，系统误差完全正相关。求 $x$ 的最佳线性无偏估计量及其不确定度。
\begin{solution}
记
\[
\bm x=\begin{pmatrix}x_1\\x_2\end{pmatrix},
\qquad
\bm 1=\begin{pmatrix}1\\1\end{pmatrix}.
\]
统计误差相互独立，而系统误差完全正相关，因此协方差矩阵为
\[
V=\begin{pmatrix}
\sigma_{s1}^2+\sigma_{c1}^2 & \sigma_{c1}\sigma_{c2}\\
\sigma_{c1}\sigma_{c2} & \sigma_{s2}^2+\sigma_{c2}^2
\end{pmatrix}.
\]
最佳线性无偏估计量，即广义最小二乘估计，为
\[
\hat x=\frac{\bm 1^TV^{-1}\bm x}{\bm 1^TV^{-1}\bm 1},
\qquad
\V(\hat x)=\frac1{\bm 1^TV^{-1}\bm 1}.
\]
为写成显式形式，令
\[
v_1=\sigma_{s1}^2+\sigma_{c1}^2,
\qquad
v_2=\sigma_{s2}^2+\sigma_{c2}^2,
\qquad
c=\sigma_{c1}\sigma_{c2}.
\]
则
\[
V=\begin{pmatrix}v_1&c\\c&v_2\end{pmatrix}.
\]
直接计算 $V^{-1}$ 可得权重
\[
w_1=\frac{v_2-c}{v_1+v_2-2c},
\qquad
w_2=\frac{v_1-c}{v_1+v_2-2c},
\qquad
w_1+w_2=1.
\]
因此
\ansmath{\hat x=w_1x_1+w_2x_2}
即
\[
\hat x=
\frac{(\sigma_{s2}^2+\sigma_{c2}^2-\sigma_{c1}\sigma_{c2})x_1
+(\sigma_{s1}^2+\sigma_{c1}^2-\sigma_{c1}\sigma_{c2})x_2}
{\sigma_{s1}^2+\sigma_{s2}^2+(\sigma_{c1}-\sigma_{c2})^2}.
\]
估计量方差为
\[
\V(\hat x)=\frac{v_1v_2-c^2}{v_1+v_2-2c}
=\frac{\sigma_{s1}^2\sigma_{s2}^2+
\sigma_{s1}^2\sigma_{c2}^2+
\sigma_{c1}^2\sigma_{s2}^2}
{\sigma_{s1}^2+\sigma_{s2}^2+(\sigma_{c1}-\sigma_{c2})^2}.
\]
所以最终结果为
\ansmath{\hat x\pm\sqrt{\V(\hat x)}}.

特别地，若两次测量的系统不确定度完全相同，即 $\sigma_{c1}=\sigma_{c2}=\sigma_c$，则
\[
w_1=\frac{\sigma_{s2}^2}{\sigma_{s1}^2+\sigma_{s2}^2},
\qquad
w_2=\frac{\sigma_{s1}^2}{\sigma_{s1}^2+\sigma_{s2}^2},
\]
并且
\[
\V(\hat x)=\frac{\sigma_{s1}^2\sigma_{s2}^2}{\sigma_{s1}^2+\sigma_{s2}^2}+\sigma_c^2.
\]
这说明共同系统误差不能通过重复测量按 $1/\sqrt n$ 消除。
\end{solution}
\end{question}